\section{Conclusion}
\label{sec:conclusion}

This paper began with a simple question about foundational growth.  When a cubical
foundation is extended by a sealed, globally acting structural layer, how much of the
resulting integration burden must remain as primitive public data?  A sealed extension is
not merely another transparent definition in a library.  It exports an interface, and that
interface must explain how the new layer acts on the active structure that the old foundation
already exposes.  The explanatory residue is public trace.

The main result is that, in the fixed sealed-extension calculus studied here, this public
structural trace has exact depth two.  Unary action trace is not enough: univalence makes
binary comparison data public, and a sealed action must respect it.  But higher
sealing-generated structural trace is not primitive.  Once the unary and binary boundary data
are present, the remaining structural obligations are open-box extension problems.  Cubical
composition, filling, and transport compute the missing faces and fillers.  Thus the exact
semantic obligation object still contains higher factors, but those factors are contractible;
the normalized opaque public signature contains no primitive structural trace fields above
binary arity.

In slogan form:
\begin{quote}
Binary trace is necessary; higher structural trace is cubically payable.
\end{quote}

This conclusion is intentionally stated for the fixed calculus \(C_{\mathrm{ext}}\).  The
calculus is not a decorative presentation of an already obvious fact.  It is the theorem
object.  Its definitions separate payload from trace, primitive fields from derived
horn-computational fields, exact realized obligations from \(\mu\)-minimal public signatures,
and arity depth from chronological support.  Those separations are what make the depth-two
claim precise.

\subsection{The theorem stack}
\label{subsec:conclusion-theorem-stack}

The proof proceeds through a stack of results, each doing a different kind of work.

First, the raw calculus gives a disciplined account of public trace.  A sealed declaration has
payload, unary action trace, binary comparison trace, and higher structural horn packages.
Only the first two trace roles are candidates for primitive structural integration cost in a
minimal public signature.  Higher structural packages are retained in the exact obligation
object, but they are marked as theorem-facing obligations whose derivedness must be justified
by cubical computation.

Second, the adequacy and normalization layer explains why the paper does not count surface
syntax directly.  Surface presentations may include redundant fields, explicit higher
coherence witnesses, or alternative notations for the same public behavior.  The trace cost
\(\mu\) is attached to normalized opaque signatures, not to arbitrary presentations.  The SMod
instance shows that this boundary is non-vacuous: a concrete surface fragment can elaborate
into the raw calculus while preserving payload, public trace, historical support, and
primitive-versus-derived status.

Third, the structural horn layer identifies the exact shape of the higher obligations.  The
important theorem is not a closed-boundary uniqueness theorem for arbitrary higher homotopy.
It is an open-box theorem for the marginal structural extension problems generated by sealing.
The lower-dimensional faces come from already exported unary and binary trace; the missing
face and interior filler are computed by the cubical Kan operations.  This is the point at
which cubicality does real work.

Fourth, the computational replacement theorem turns semantic fillability into public-trace
elimination.  If a presentation contains an explicit higher structural field, the normalizer
replaces it by a cubical term synthesized from its boundary data.  After this replacement the
field is absent from the \(\mu\)-minimal primitive signature.  This is the reason that
``derived'' is not merely a syntactic label in the intended theorem: it is backed by a
replacement operation and a public-cost invariance statement.

Fifth, the exact depth theorem combines the cubical upper bound with the univalent lower
bound.  The upper bound says that all structural obligations of arity greater than two are
contractible over the binary boundary and therefore primitive-publicly eliminable.  The lower
bound uses the two-point swap obstruction: unary action data alone cannot in general decide
whether a sealed action respects the public comparison induced by a nontrivial equivalence.
Together these results give exact structural integration depth two for \(C_{\mathrm{ext}}\):
\[
  \mathsf{Real}_k(X) \simeq \mathsf{Real}_2(X)
  \qquad (k \geq 2),
\]
while depth one fails in general.

Finally, the counting theorems show what happens after the semantic depth theorem is combined
with additional locality and coverage hypotheses.  Depth as arity does not by itself imply a
chronological recurrence: a binary comparison can still refer to remote history.  The
recent-history theorem adds a factorization-complete export condition under which remote
comparisons factor through the two most recent normalized signatures.  Full coupling then
turns the semantic window into arithmetic:
\[
  \mu_{n+1}
  =
  \mu_n + \mu_{n-1} + \kappa_n + \kappa_{n-1}.
\]
The Fibonacci endpoint is the constant-payload specialization of this recurrence, not the
central theorem of the paper.

\subsection{What the result does not say}
\label{subsec:conclusion-nonclaims}

Several nonclaims are as important as the theorem itself.

The result does not say that higher paths disappear from cubical type theory.  Cubical
foundations are rich precisely because higher structure is available and computationally
meaningful.  The theorem concerns the primitive public trace forced by sealing a globally
acting structural extension.  User-supplied higher payload, internal higher homotopy, and
ordinary higher-dimensional constructions remain part of the ambient theory.

The result also does not say that every higher filler space in homotopy type theory is
contractible.  The contractibility theorem is about theorem-facing structural open-box
extension fibers generated by the sealed-extension calculus.  These fibers include boundary,
side, lid, and agreement data arranged so that cubical composition and filling provide the
canonical extension.  Arbitrary closed higher-dimensional choices are not being collapsed.

Nor does the recurrence claim describe ordinary library growth.  Most library development is
sparse.  A new transparent definition typically touches only a small part of the active
interface, and even sealed extensions may interact with a sparse footprint.  The full-coupling
recurrence applies only under an explicit regime: depth-two public trace, two-layer
chronological factorization, total active-basis coverage, and the bundled per-site trace
convention.  Without these hypotheses the correct formula is the sparse footprint formula,
not the full Fibonacci-style envelope.

Finally, the paper does not prove an unrestricted adequacy theorem for every cubical
foundation or every surface language.  It proves a fixed-calculus theorem and isolates a
transfer interface.  A new system inherits the theorem only after supplying a faithful
interpretation that preserves semantic admissibility, primitive-versus-derived trace status,
historical support, horn-computational replacement, and public trace cardinality.

\subsection{Why the depth-two boundary is natural}
\label{subsec:conclusion-why-depth-two}

The depth-two boundary is not an arbitrary compromise between unary simplicity and an
infinite coherence tower.  It reflects two different sources of structure.

The lower bound comes from public comparison.  In a univalent setting, equivalences are not
invisible bookkeeping; they determine paths and therefore become part of the public structure
that a sealed action must respect.  Objectwise action cannot see all of this structure.  A
sealed extension that only records unary behavior may agree on old generators while failing to
respect a public equivalence between them.  Binary trace is therefore forced by the semantics
of the foundation.

The upper bound comes from cubical computation.  Once the binary boundary is present, the
higher structural marginal obligations have the form for which the cubical core was designed:
open boxes with compatible partial boundaries.  The higher obligation is not a new independent
public generator; it is the result of solving an extension problem whose boundary is already
public.  Composition and filling are exactly the operations that solve such problems.

This is why the theorem has the shape it does.  A non-cubical or extensional setting with UIP
can collapse comparison data further, yielding a depth-one contrast.  A weak system without
substitution-stable filling may fail the upper-bound mechanism.  A surface language without an
adequate sealed-interface normalization may fail to transfer the theorem.  But in the fixed
cubical calculus, univalence supplies the binary lower bound and Kan filling supplies the
higher upper bound.

\subsection{The role of mechanization}
\label{subsec:conclusion-mechanization}

The accompanying formalization is meant to support exactly this theorem stack.  It checks the
fixed raw calculus, the structural open-box and horn-computation layer, the derivedness and
trace-accounting discipline, the SMod adequacy instance, the lower-bound witness, and the
counting arithmetic.  The theorem map and audit scripts are included to make the boundary of
formal evidence explicit.

The mechanization is valuable for two reasons.  First, it guards against the easiest mistake in
this project: proving derivedness by tagging a field as derived.  The important replacement
result must show that a higher structural field elaborates to a cubical expression and then
vanishes from the minimal primitive signature.  Second, the formal development forces the
open-box theorem to state its boundary, side, lid, filler, and agreement data explicitly,
rather than hiding the hard semantic step inside a singleton abstraction.

At the same time, mechanization does not remove the need for the paper's semantic
formulation.  A checked theorem proves the statement encoded by the checked definitions.  The
paper proofs explain why those definitions are the intended model of sealed structural trace,
which claims are internal to \(C_{\mathrm{ext}}\), and which claims require external adequacy
hypotheses.  This combination---a fixed checked theorem plus an explicit transfer
criterion---is the intended standard of evidence.

\subsection{Open directions}
\label{subsec:conclusion-open-directions}

The most important next direction is adequacy transfer.  The abstract horn-computational
interface gives a precise target: a candidate foundation or surface language should provide a
sealed-interface interpretation preserving admissibility, public normalization, historical
support, primitive-versus-derived classification, and trace cardinality.  Proving this for
additional modal fragments, universe extensions, or module systems would test how robust the
fixed-calculus theorem is.

A second direction is to refine the open-box semantics.  The present calculus isolates the
structural extension fibers needed for the depth theorem.  Future work could compare this
presentation with other cubical models, other notions of partial element and composition, and
other forms of opacity.  Such comparisons would clarify which parts of the theorem depend on
CCHM-style operations and which parts follow from more abstract Kan or cofibration structure.

A third direction is to study weaker and stronger locality regimes.  The paper separates arity
depth from chronological support, but the two-layer Markov blanket is only one possible
factorization principle.  Other extension disciplines may have one-layer, two-layer,
finite-depth, or sparse nonuniform historical windows.  The sparse footprint formulas suggest
that the right general theory is not a single recurrence, but a family of trace-cost laws
parameterized by support and coupling data.

A fourth direction is to sharpen the relationship between active-basis coverage and
canonicity.  In this paper, full coupling is an explicit regime used by the recurrence theorem.
A deeper semantic analysis might identify classes of foundational-core extensions for which
advertised total global action really does force per-site primitive coverage, modulo
normalization and derived-field elimination.  That would turn part of the current counting
hypothesis into a theorem for those systems.

\subsection{Final perspective}
\label{subsec:conclusion-final-perspective}

The guiding intuition of the paper is that opacity changes what a foundation owes its users.
A transparent construction may rely on its context without publishing a complete global
interaction contract.  A sealed foundational layer cannot.  Once exported, it must carry enough
public trace for future code to use it without reopening its implementation.

Cubical structure makes that contract surprisingly economical.  The first nontrivial public
obligation is binary, because a univalent foundation exposes comparisons as structure.  But the
apparent higher tower is not an infinite list of new primitive clauses.  It is a family of
open-box extension problems whose solutions are computed by the cubical core.  The public
signature therefore stabilizes at binary trace even though the exact semantic object continues
to remember the higher contractible witnesses.

This is the sense in which coherence depth two is both necessary and enough.  It is necessary
because unary action cannot see public equivalence.  It is enough because higher structural
coherence, once correctly formulated as marginal cubical extension, is not a new primitive
public debt.  The broader conjecture is that this pattern is not an accident of the chosen
calculus, but a recurring feature of cubical foundations with stable sealing, univalence, and
Kan computation.  The fixed theorem proved here is one precise instance of that pattern, and a
framework for testing the others.
